% All figures at publication quality, one per page, with their legends.
% Filled by R/08_assemble_reports.R and compiled with XeLaTeX from report/.
\documentclass[10pt]{article}
\usepackage[letterpaper,margin=0.8in]{geometry}
\usepackage{fontspec}
\setmainfont{Liberation Sans}
\usepackage{graphicx,xcolor}
\usepackage[hidelinks]{hyperref}
\setlength{\parindent}{0pt}
\newcommand{\figpage}[3]{\clearpage\begin{center}{\large\bfseries #1}\par\vspace{8pt}\includegraphics[width=#2\textwidth]{#3}\end{center}\vspace{6pt}}
\begin{document}
{\LARGE\bfseries Experiment 2: β-galactosidase kinetics}\par\vspace{4pt}
{\large Figures (main and supplementary) at publication quality}\par\vspace{10pt}
Almokhtar Aljarodi and Abdulrahman Khardali, Methods in Molecular and Cellular Biology (B241), KAUST. Compiled 10 October 2026.\par\vspace{10pt}
Main figures are drawn from the reference dataset (course-provided measurements used for the Results section at the instructor's direction); supplementary figures hold the experimental dataset (our own printouts), diagnostics and sensitivity analyses. Conventions: filled symbol, observation used for \textit{V}\textsubscript{0}; open symbol, observation retained but not used; solid line, fitted initial window; dashed segment from the origin, 0–2 min secant through the assumed origin (no reading exists at t = 0); dotted line, extrapolation; ↑ marks a lower-bound \textit{V}\textsubscript{0}. Each condition was measured once, so no error bars are shown. Full legends are in the course report and in the journal-style version; the vector PDF of every figure is in \texttt{figures/}.\par\vspace{10pt}
Contents: Figure 1 (ONP spectrum and standard curve), Figure 2 (enzyme dilution), Figure 3 (substrate saturation, no inhibitor), Figure 4 (IPTG inhibition and Lineweaver–Burk plots), graphical abstract, Figures S1–S4, reference diagram for Question 1.

\figpage{Figure 1. ONP absorbs maximally at 415 nm and obeys the Beer–Lambert law to 0.43 mM}{1}{../figures/figure1_onp_calibration.pdf}
(A) Absorbance spectrum of 0.333 mM ONP in the alkaline assay matrix; dashed line, maximum at 415 nm; shaded band, within 2\% of the maximum. (B) Standard curve at 415 nm; filled points and line, least-squares calibration over the shaded range, \textit{A}\textsubscript{415} = 5.222[ONP] − 0.022, \textit{R}\textsuperscript{2} = 0.9988, ε = 5.22 ± 0.10 mM\textsuperscript{−1} cm\textsuperscript{−1}; ×, zero-concentration point (not fitted); open point, excluded 0.533 mM standard (17\% above the line). Reference dataset, Part I.

\figpage{Figure 2. Initial velocity rose approximately in proportion to the amount of β-galactosidase}{1}{../figures/figure2_enzyme_dilution.pdf}
(A) Progress curves for 2.20 mM ONPG with 1.00 to 0.15 relative enzyme (tubes A–E); lines, least-squares fits to the initial linear window; filled points used, open points excluded, dotted extrapolation. (B) \textit{V}\textsubscript{0} in µmol ONP min\textsuperscript{−1} per 12.0 mL assay against relative enzyme; solid line through the origin, slope 0.481, \textit{R}\textsuperscript{2} = 0.972; dotted, free-intercept fit (intercept -0.0006, p = 0.98). Relative rates 1.00, 0.77, 0.60, 0.18, 0.17 for relative enzyme 1.00, 0.75, 0.50, 0.25, 0.15. Reference dataset, Part II.

\figpage{Figure 3. β-Galactosidase approached saturation with ONPG}{1}{../figures/figure3_substrate_saturation.pdf}
(A) Progress curves for 0.2 mL enzyme with 2.20 to 0.069 mM ONPG (tubes A–F, no inhibitor); tube A linear (regression); in B–F (†) the rate fell within the first intervals, so \textit{V}\textsubscript{0} is the 0–2 min secant (dashed) through the assumed origin (×), a lower bound; open points not used. (B) \textit{V}\textsubscript{0} against [ONPG]; solid, nonlinear fit (\textit{K}\textsubscript{m} = 0.36 ± 0.13 mM, \textit{V}\textsubscript{max} = 0.494 ± 0.059 µmol min\textsuperscript{−1}); dashed, hyperbola from the Lineweaver–Burk line (\textit{K}\textsubscript{m} = 0.31 mM, \textit{V}\textsubscript{max} = 0.477); ↑ lower bound. Reference dataset, Part III.

\figpage{Figure 4. IPTG raised the apparent \textit{K}\textsubscript{m} of β-galactosidase for ONPG}{1}{../figures/figure4_iptg_inhibition.pdf}
(A) \textit{V}\textsubscript{0} against [ONPG] without inhibitor (circles) and with 0.458 mM IPTG (triangles); dashed, hyperbolae from the Lineweaver–Burk lines; solid, one competitive model fitted to both series (shared \textit{V}\textsubscript{max} = 0.53 µmol min\textsuperscript{−1}, \textit{K}\textsubscript{m} = 0.44 mM, \textit{K}\textsubscript{i} = 0.107 ± 0.042 mM; the shared \textit{V}\textsubscript{max} is a model assumption); ↑ lower bound. (B) Lineweaver–Burk plot: no inhibitor 1/\textit{V}\textsubscript{0} = 0.66(1/[S]) + 2.10 (\textit{K}\textsubscript{m} 0.31 mM, \textit{V}\textsubscript{max} 0.477); + IPTG 1/\textit{V}\textsubscript{0} = 6.24(1/[S]) + 3.55 (\textit{K}\textsubscript{m,app} 1.76 mM, \textit{V}\textsubscript{max,app} 0.282); \textit{K}\textsubscript{i} = 0.100 mM by the competitive relation. Reference dataset, Part III.

\figpage{Graphical abstract}{0.95}{../figures/graphical_abstract.pdf}
Reaction scheme, assay logic and the two saturation curves of the reference dataset with the headline values; the competitive model's shared \textit{V}\textsubscript{max} is a model assumption. \textit{K}\textsubscript{i} is given as the range obtained across initial-velocity estimators.

\figpage{Figure S1. Progress curves of the two datasets on common axes}{0.9}{../figures/supplementary/figureS1_experimental_vs_reference_curves.pdf}
Reference dataset (left) and experimental dataset (right) for the three experiments; solid segments, regression windows; dashed, 0–2 min secants; open, not used for \textit{V}\textsubscript{0}. The experimental inhibited tube B (open) is excluded from fits; inhibited tube E has no data; tube F was not performed in the experimental sessions.

\figpage{Figure S2. Increment diagnostics behind every \textit{V}\textsubscript{0} decision}{1}{../figures/supplementary/figureS2_increment_diagnostics.pdf}
Each 2-min increment divided by the increment predicted by the fitted \textit{V}\textsubscript{0}; shaded band, acceptance window of the rule (±20\% or ±0.010 A); open, intervals not used. For secant-based curves the reference rate is the 0–2 min secant, so later intervals fall below 1.

\figpage{Figure S3. Sensitivity of the kinetic parameters to the initial-velocity estimator}{1}{../figures/supplementary/figureS3_estimator_sensitivity.pdf}
\textit{K}\textsubscript{m}, \textit{K}\textsubscript{m,app} and \textit{K}\textsubscript{i} of the reference dataset under five estimators (A, all five points; R, regression only; S, 0–2 min secant only; O, rule with offset-corrected secants; P, the rule used). Cyan lines, values for purified \textit{E. coli} β-galactosidase (TES buffer, pH 7.0, 25 °C).

\figpage{Figure S4. Kinetic analysis of the experimental dataset}{1}{../figures/supplementary/figureS4_experimental_kinetics.pdf}
(A) Michaelis–Menten representation of our Part III data (nonlinear fits; open, excluded inhibited tube B). (B) Ratio of our \textit{V}\textsubscript{0} to the reference \textit{V}\textsubscript{0} for every tube measured in both datasets.

\figpage{Reference diagram for Question 1 (schematic, to be redrawn by hand)}{0.85}{../figures/supplementary/figureQ1_michaelis_menten_sketch.pdf}
Michaelis–Menten curve to [S] = 4 \textit{K}\textsubscript{m} with the asymptote V = \textit{V}\textsubscript{max}, the line [S] = \textit{K}\textsubscript{m} and the tangent at the origin (slope \textit{V}\textsubscript{max}/\textit{K}\textsubscript{m}); V(4\textit{K}\textsubscript{m}) = 0.8 \textit{V}\textsubscript{max}. This computer-drawn diagram is a reference, not the hand-drawn sketch the course requires.
\end{document}
